\documentclass[12pt,a4paper]{article}
\usepackage[margin=1in]{geometry}
\usepackage[T1]{fontenc}
\usepackage[utf8]{inputenc}
\usepackage{graphicx}
\usepackage{lmodern}
\usepackage{setspace}
\usepackage{enumitem}
\usepackage{amssymb}
\usepackage{needspace}

\setstretch{1.1}
\setlength{\parindent}{0pt}
\setlength{\parskip}{0.6em}

\newcommand{\boxmark}{\(\square\)}
\newcommand{\lineanswer}{\vspace{2.5em}\hrulefill}
\newcommand{\shortline}{\underline{\hspace{3cm}}}
\newcommand{\yesno}{\boxmark\ Yes \hspace{2em} \boxmark\ No}

\newcommand{\responseline}[5]{%
    \begin{flushleft}
    \boxmark\ #1 \hspace{1.2em}
    \boxmark\ #2 \hspace{1.2em}
    \boxmark\ #3 \hspace{1.2em}
    \boxmark\ #4 \hspace{1.2em}
    \boxmark\ #5
    \end{flushleft}
}

\newcommand{\fivepointscale}[1]{%
    \Needspace{4\baselineskip}%
    \textbf{#1}
    
    \responseline{Strongly disagree}{Disagree}{Neither agree nor disagree}{Agree}{Strongly agree}
}

\newcommand{\familiarityscale}[1]{%
    \Needspace{4\baselineskip}%
    \textbf{#1}
    
    \responseline{None}{Limited}{Moderate}{Good}{High}
}

\newcommand{\difficultyscale}[1]{%
    \Needspace{4\baselineskip}%
    \textbf{#1}
    
    \responseline{Very easy}{Easy}{Neutral}{Difficult}{Very difficult}
}

\newcommand{\confidenceoptional}{%
    \textit{Optional confidence rating:} \boxmark\ 1 \hspace{1em} \boxmark\ 2 \hspace{1em} \boxmark\ 3 \hspace{1em} \boxmark\ 4 \hspace{1em} \boxmark\ 5
}

\begin{document}

\begin{flushleft}
    \includegraphics[width=0.28\textwidth]{Logo.png}
\end{flushleft}

\begin{center}
    {\LARGE \textbf{Study Questionnaire}}\\[0.4em]
    {\large Creating Novel Virtual Reality Interfaces for Teaching Computer Architecture}
\end{center}

\textbf{Participant ID:} \shortline \hfill \textbf{Date:} \shortline

\section*{Section A: Screening and Background}

\textbf{1. Are you 18 or older?}

\yesno

\textbf{2. Have you previously studied computer architecture or computer organisation?}

\yesno

\familiarityscale{3. How familiar are you with virtual reality?}

\textbf{4. Have you used a VR headset before?}

\yesno

\fivepointscale{5. Before this session, how confident are you in your understanding of computer architecture concepts?}

\difficultyscale{6. How difficult do you usually find topics such as instruction flow, pipelining, stack behaviour, and memory hierarchy?}

\textbf{7. Which of the following topics have you previously studied? Tick all that apply.}

\boxmark\ Instruction flow \hspace{1em}
\boxmark\ Pipelining \hspace{1em}
\boxmark\ Stack \hspace{1em}
\boxmark\ Memory hierarchy \hspace{1em}
\boxmark\ Assembly \hspace{1em}
\boxmark\ None

\section*{Section B: System Usability}

\textit{For Questions 8--14, please select one response for each statement.}

\fivepointscale{8. The VR system was easy to learn.}
\fivepointscale{9. I could navigate the environment without difficulty.}
\fivepointscale{10. The controls felt intuitive.}
\fivepointscale{11. The visual presentation made the concepts easier to follow.}
\fivepointscale{12. I could focus on the learning task without being distracted by the interface.}
\fivepointscale{13. I would feel confident using this system again.}
\fivepointscale{14. I experienced discomfort, motion sickness, or visual fatigue while using the system.}

\section*{Section C: System Engagement}

\textit{For Questions 15--19, please select one response for each statement.}

\fivepointscale{15. The experience held my attention.}
\fivepointscale{16. I felt actively involved in the learning process.}
\fivepointscale{17. The system made the topic feel more interesting.}
\fivepointscale{18. I was motivated to explore the concepts further.}
\fivepointscale{19. The interactive format improved my focus compared with standard learning materials.}

\section*{Section D: Perceived Learning}

\textit{For Questions 20--24, please select one response for each statement.}

\fivepointscale{20. The VR experience helped me understand the topic better.}
\fivepointscale{21. The visualisations helped me build a mental model of what was happening.}
\fivepointscale{22. The system made abstract processes easier to understand.}
\fivepointscale{23. I would prefer this as a supplement to traditional teaching materials.}
\fivepointscale{24. I believe I learned something useful from this session.}

\section*{Section E: Conceptual Questions}

\textit{For each question, tick one answer. A confidence rating can be included if you want extra analysis, but I recommend keeping it optional rather than mandatory unless confidence is a core study outcome.}

\textbf{25. Which register stores the address of the next instruction?}

\boxmark\ Stack Pointer \hspace{2em}
\boxmark\ Program Counter \hspace{2em}
\boxmark\ Instruction Register \hspace{2em}
\boxmark\ Base Register

\confidenceoptional

\textbf{26. Which stage of the instruction cycle reads an instruction from memory?}

\boxmark\ Decode \hspace{2em}
\boxmark\ Fetch \hspace{2em}
\boxmark\ Execute \hspace{2em}
\boxmark\ Write-back

\confidenceoptional

\textbf{27. During instruction execution, the instruction register primarily stores:}

\boxmark\ Program variables \hspace{1em}
\boxmark\ The currently executing instruction \hspace{1em}
\boxmark\ The stack pointer value \hspace{1em}
\boxmark\ Cache addresses

\confidenceoptional

\textbf{28. What happens to the program counter after an instruction is fetched?}

\boxmark\ It resets to zero \hspace{1em}
\boxmark\ It increments to the next instruction address \hspace{1em}
\boxmark\ It copies into the stack pointer \hspace{1em}
\boxmark\ It stores the result of the ALU

\confidenceoptional

\textbf{29. What is the main purpose of the stack during function calls?}

\boxmark\ Store pipeline stages \hspace{1em}
\boxmark\ Store return addresses and local variables \hspace{1em}
\boxmark\ Store instruction cache data \hspace{1em}
\boxmark\ Store CPU clock values

\confidenceoptional

\textbf{30. Which register tracks the top of the stack?}

\boxmark\ PC \hspace{2em}
\boxmark\ SP \hspace{2em}
\boxmark\ IR \hspace{2em}
\boxmark\ ALU

\confidenceoptional

\textbf{31. When a function returns, which of the following usually happens?}

\boxmark\ Stack pointer moves back to previous position \hspace{1em}
\boxmark\ Program counter resets to zero \hspace{1em}
\boxmark\ Cache is cleared \hspace{1em}
\boxmark\ Registers are deleted

\confidenceoptional

\textbf{32. Which event normally pushes data onto the stack?}

\boxmark\ Memory read \hspace{2em}
\boxmark\ Function call \hspace{2em}
\boxmark\ Pipeline stall \hspace{2em}
\boxmark\ Cache miss

\confidenceoptional

\textbf{33. Why do CPUs use pipelining?}

\boxmark\ To reduce instruction memory \hspace{1em}
\boxmark\ To increase instruction throughput \hspace{1em}
\boxmark\ To eliminate branching \hspace{1em}
\boxmark\ To reduce stack usage

\confidenceoptional

\textbf{34. Which of the following best describes a pipeline stall?}

\boxmark\ A new instruction enters the pipeline \hspace{1em}
\boxmark\ The pipeline pauses due to a dependency or hazard \hspace{1em}
\boxmark\ The stack pointer changes \hspace{1em}
\boxmark\ The ALU executes multiple instructions

\confidenceoptional

\textbf{35. A data hazard occurs when:}

\boxmark\ An instruction depends on the result of a previous instruction \hspace{1em}
\boxmark\ Memory latency increases \hspace{1em}
\boxmark\ Branch prediction fails \hspace{1em}
\boxmark\ Cache memory overflows

\confidenceoptional

\textbf{36. Which pipeline stage usually performs arithmetic operations?}

\boxmark\ Fetch \hspace{2em}
\boxmark\ Decode \hspace{2em}
\boxmark\ Execute \hspace{2em}
\boxmark\ Write-back

\confidenceoptional

\textbf{37. Branch instructions primarily affect which component?}

\boxmark\ Cache memory \hspace{1em}
\boxmark\ Program counter \hspace{1em}
\boxmark\ Stack pointer \hspace{1em}
\boxmark\ Instruction register

\confidenceoptional

\textbf{38. If a branch prediction is incorrect:}

\boxmark\ The CPU must flush incorrect instructions in the pipeline \hspace{1em}
\boxmark\ The stack pointer resets \hspace{1em}
\boxmark\ The instruction register deletes memory \hspace{1em}
\boxmark\ The ALU halts

\confidenceoptional

\textbf{39. Why are caches used in modern CPUs?}

\boxmark\ To increase instruction size \hspace{1em}
\boxmark\ To store program code permanently \hspace{1em}
\boxmark\ To reduce average memory access time \hspace{1em}
\boxmark\ To remove the need for RAM

\confidenceoptional

\textbf{40. Which memory type is fastest?}

\boxmark\ Hard disk \hspace{2em}
\boxmark\ RAM \hspace{2em}
\boxmark\ Cache \hspace{2em}
\boxmark\ Virtual memory

\confidenceoptional

\textbf{41. Which memory component is typically largest but slowest?}

\boxmark\ Cache \hspace{2em}
\boxmark\ Registers \hspace{2em}
\boxmark\ Main memory \hspace{2em}
\boxmark\ Secondary storage

\confidenceoptional

\textbf{42. What does this instruction do? \texttt{add x1, x2, x3}}

\boxmark\ Adds \texttt{x2} and \texttt{x3} and stores the result in \texttt{x1}\\
\boxmark\ Adds \texttt{x1} and \texttt{x2} and stores the result in \texttt{x3}\\
\boxmark\ Moves \texttt{x2} into \texttt{x3}\\
\boxmark\ Multiplies registers

\confidenceoptional

\textbf{43. What does this instruction do? \texttt{addi x5, x0, 10}}

\boxmark\ Stores 10 into \texttt{x5} \hspace{1em}
\boxmark\ Clears \texttt{x5} \hspace{1em}
\boxmark\ Moves \texttt{x5} into \texttt{x0} \hspace{1em}
\boxmark\ Adds two registers

\confidenceoptional

\textbf{44. Which instruction type loads data from memory into a register?}

\boxmark\ ADD \hspace{2em}
\boxmark\ STORE \hspace{2em}
\boxmark\ LOAD \hspace{2em}
\boxmark\ JUMP

\confidenceoptional

\textbf{45. Consider this sequence: \texttt{lw x1, 0(x2)} followed by \texttt{add x3, x1, x4}. Why might a pipeline stall occur?}

\boxmark\ Because the second instruction depends on \texttt{x1}\\
\boxmark\ Because \texttt{x2} is unused\\
\boxmark\ Because the program counter resets\\
\boxmark\ Because memory is cleared

\confidenceoptional

\textbf{46. Which register usually holds the return address after a function call in RISC-V?}

\boxmark\ \texttt{x0} \hspace{2em}
\boxmark\ \texttt{ra} \hspace{2em}
\boxmark\ \texttt{sp} \hspace{2em}
\boxmark\ \texttt{gp}

\confidenceoptional

\textbf{47. What is the role of register \texttt{x0} in RISC-V?}

\boxmark\ Stores the program counter \hspace{1em}
\boxmark\ Always contains the value 0 \hspace{1em}
\boxmark\ Stores return addresses \hspace{1em}
\boxmark\ Tracks the stack pointer

\confidenceoptional

\textbf{48. What does this instruction do? \texttt{lw x5, 0(x10)}}

\boxmark\ Stores \texttt{x5} into memory \hspace{1em}
\boxmark\ Loads a word from memory into \texttt{x5} \hspace{1em}
\boxmark\ Moves \texttt{x10} into \texttt{x5} \hspace{1em}
\boxmark\ Clears memory

\confidenceoptional

\textbf{49. What does this instruction do? \texttt{sw x5, 4(x10)}}

\boxmark\ Value from \texttt{x5} is stored in memory \hspace{1em}
\boxmark\ Value from memory is loaded into \texttt{x5} \hspace{1em}
\boxmark\ \texttt{x10} is overwritten \hspace{1em}
\boxmark\ Stack pointer moves

\confidenceoptional

\textbf{50. What is this instruction commonly used for? \texttt{addi sp, sp, -16}}

\boxmark\ Allocate space on the stack \hspace{1em}
\boxmark\ Reset stack pointer \hspace{1em}
\boxmark\ Jump to a function \hspace{1em}
\boxmark\ Clear registers

\confidenceoptional

\textbf{51. What is this instruction commonly used for? \texttt{addi sp, sp, 16}}

\boxmark\ Allocates stack memory \hspace{1em}
\boxmark\ Frees stack memory after a function returns \hspace{1em}
\boxmark\ Loads memory \hspace{1em}
\boxmark\ Changes the program counter

\confidenceoptional

\textbf{52. What does this instruction do? \texttt{beq x1, x2, label}}

\boxmark\ Jump if \texttt{x1 > x2} \hspace{1em}
\boxmark\ Jump if \texttt{x1 == x2} \hspace{1em}
\boxmark\ Jump if \texttt{x1 < x2} \hspace{1em}
\boxmark\ Always jump

\confidenceoptional

\textbf{53. What does this instruction do? \texttt{jal ra, func}}

\boxmark\ Calls a function and stores the return address in \texttt{ra}\\
\boxmark\ Jumps without saving the return address\\
\boxmark\ Loads memory\\
\boxmark\ Writes to the stack

\confidenceoptional

\textbf{54. Put the instruction stages in the correct order.}

\boxmark\ Execute, Fetch, Decode, Write-back\\
\boxmark\ Fetch, Decode, Execute, Write-back\\
\boxmark\ Decode, Fetch, Execute, Write-back\\
\boxmark\ Fetch, Execute, Decode, Write-back

\confidenceoptional

\section*{Section F: Open-Ended Feedback}

\textbf{55. What part of the experience helped your understanding most?}
\lineanswer

\textbf{56. What part was confusing or ineffective?}
\lineanswer

\textbf{57. What would you improve in the system?}
\lineanswer

\textbf{58. Were there any moments where the VR interface got in the way of learning?}
\lineanswer

\textbf{59. Would you use this again as a learning aid? Why or why not?}

\vspace{4em}
\hrulefill

\end{document}
